\section{A weighted Hessian inequality for the Weyl blocks}
\label{sec:hessian}

We continue the case in which neither Weyl block vanishes identically.
The preceding argument gives
\(\mathbf E v^2=\mathbf E b^2\). This equality will enter the coupling
in Section~\ref{sec:coupling}; the weighted Hessian inequality proved here
does not itself require it. Our task is to extract a first-derivative
inequality from the Bianchi-constrained second derivatives. Unequal
weights on the two blocks will create a useful interaction, but their
derivatives have costs that must be retained exactly.
Retain the notation
\(Z=\nabla T\), \(p=\nabla v\), \(q=\nabla b\), and
\(e=|Z|^2/6\).
The argument is stated for $T$; exchanging the two orientations gives
its companion for $U$. All expectations use normalized Riemannian volume.
The constrained decomposition and weighted square come from
\cite[Lemmas~5.1--5.2 and Proposition~5.3]{OpenAI2026P}. We prove the
identities in full and retain their nonnegative form at the sectional
boundary. The square will first be proved for an arbitrary global positive
Lipschitz weight. We choose the two fixed weights when their derivative
costs enter.

\subsection{The constrained Hessian and its curvature operator}

Retain $F=\operatorname{Sym}_0(\Lambda^+)$ and the Bianchi gradient
bundle $G\subset T^*M\otimes F$ from Section~\ref{sec:preliminaries}.
Their types are $(2,0)$ and $(5/2,1/2)$. The irreducible gradient
constraint and the refined Kato coefficient used below are established
in Gursky--LeBrun~\cite[Lemma~4, Equations~(10)--(12)]{GurskyLeBrun1999};
Calderbank--Gauduchon--Herzlich~\cite{CalderbankGauduchonHerzlich2000}
place them in the general generalized-gradient framework. We use the
parallel projector $P_G=(6-C)/10$ and the compression identities already
proved in Proposition~\ref{prop:weyl-derivative-equations}. They are
pointwise tensor identities, including at repeated eigenvalues, and do
not require a smoothly chosen eigenbasis.

\begin{lemma}[The two Hessian components]
\label{lem:hessian-components}
Let \(\mathcal H=T^*M\otimes G\), and write
\(X=\nabla Z\), so \(X_{ij}=\nabla_i\nabla_jT\).
Then
\[
 X=Y+D_0,
 \qquad
 Y\text{ has type }(3,1),\qquad
 D_0\text{ has type }(2,0).
\]
If \(\operatorname{Tr}:\mathcal H\to F\) is contraction of the
two derivative slots and \(\operatorname{Tr}^*\) its adjoint, then
\[
 \operatorname{Tr}\operatorname{Tr}^*=\frac{12}{5}\mathrm{Id},
 \qquad
 D_0=\frac5{12}\operatorname{Tr}^*\Delta T,
 \qquad
 |D_0|^2=\frac5{12}|\Delta T|^2.
\]
All the type projections in this statement are parallel.
\end{lemma}

\begin{proof}
The product rule gives four inequivalent, multiplicity-one summands
\[
 \mathcal H=(3,1)\oplus(3,0)\oplus(2,1)\oplus(2,0).
\]
Skewing the derivative slots takes values in
\[
 \Lambda^2T^*M\otimes F=(3,0)\oplus(2,0)\oplus(1,0)\oplus(2,1).
\]
For the actual Hessian $X=\nabla\nabla T$, its skew difference is
$[\nabla_i,\nabla_j]T=F_{ij}T$, where $F_{ij}$ is the curvature of the
bundle $F$. It is essential that the tensor on which this curvature
acts is the same Weyl block $T$ that occurs in $F_{ij}$. The minus
rotation algebra acts trivially on $F$, so there is no minus component.
The dependence on $T$ is an equivariant linear term and an equivariant
term in its symmetric square. The latter has plus spins $4,2,0$;
its only common type with $\Lambda^+\otimes F=(3,0)\oplus(2,0)\oplus(1,0)$
is $(2,0)$. The linear term has that type as well. Thus the skew
commutator has neither a $(3,0)$ nor a $(2,1)$ component. This conclusion
is specific to the Hessian of $T$, not to the derivative of an arbitrary
section of $G$.

Skewing is injective on each of those two domain summands.  The
\((3,0)\) summand is entirely skew: its stretched plus coupling makes the
two derivative plus spin halves symmetric, while their minus coupling
to spin zero is alternating.  To treat \((2,1)\), suppose that a tensor
of this type is symmetric in its full derivative slots.  Its derivative
minus pair has spin \(1\) and is symmetric, so its derivative plus pair
must also be symmetric and have spin \(1\).  The inner Bianchi constraint
makes the coupling operator \(\sum I_s^{(j)}\Sigma_s\) equal to \(-4\).
To see that symmetry gives the same value for the outer coupling,
let $S$ interchange the derivative slots. Then
$C^{(i)}=SC^{(j)}S$; hence $SX=X$ and $C^{(j)}X=-4X$ imply
$C^{(i)}X=-4X$.  The total plus Casimir would therefore be
\[
 -8-24+2(-4-4)=-48,
\]
whereas total spin \(2\) requires \(-24\).  This contradiction proves
injectivity.  Moreover, the two images lie in distinct target isotypic
components, so they cannot cancel one another.  The corresponding
components of \(X\) must vanish separately.

For the trace calculation, let
\(K:F\to T^*M\otimes T^*M\otimes F\) be
\(K(S)_{ij}=\delta_{ij}S\).  The adjoint of the restricted trace is
\(\operatorname{Tr}^*=P_{G}^{(j)}K\), where the projector acts
on the inner derivative and matrix slots.  Now
\[
 \operatorname{Tr}K=4\mathrm{Id},
 \qquad
 \operatorname{Tr}\bigl(C^{(j)}K(S)\bigr)
 =\sum_{i,s}(I_s)_{ii}\Sigma_sS=0.
\]
Thus \(\operatorname{Tr}\operatorname{Tr}^*=12\mathrm{Id}/5\).
By equivariance, trace vanishes on \((3,1)\) and is supported on the
unique \((2,0)\) summand.  The displayed nonzero composition makes its
restriction to that five-dimensional summand an isomorphism onto
\(F\).  It therefore determines the whole of \(D_0\), including
its skew part, and gives both the formula and the norm identity.
One can also see explicitly that the trace component retains the
curvature commutator. At a point, its adjoint is
\[
 (\operatorname{Tr}^*S)_{ij}
 =\frac35\delta_{ij}S+\frac1{10}\sum_s I_{s,ij}\Sigma_sS.
\]
Hence
\[
 (D_0)_{ij}-(D_0)_{ji}
 =\frac1{12}\sum_s I_{s,ij}\Sigma_s\Delta T
 =\frac12\sum_s(1+t_s)I_{s,ij}\Sigma_sT.
\]
For the last equality use $\Delta T=6(T-N)$ and, in a pointwise
diagonal frame, $\Sigma_sN=-t_s\Sigma_sT$. The invariant version is
$\Sigma_rN=-\sum_sT_{rs}\Sigma_sT$; no eigenframe has been
differentiated. Thus the displayed skew difference is exactly the
bundle-curvature action.

Finally, the type projections are invariant expressions in the combined
Casimirs, and trace is a metric contraction.  They are parallel and do
not depend on any curvature eigenframe.  The spin labels describe tensor
representations of the oriented frame group; no spin structure is used.
\end{proof}

The derivative-slot operators $T_D,U_D$ were defined in the weighted
variance proof, and $B_F(T)$ acts on the matrix slot. Compress all three
orthogonally to $\mathcal H=T^*M\otimes G$. The useful combination is
visible before the quadratic form is chosen. Let
$\Sigma_s^{\mathrm{tot}}$ be the total plus generator on $\mathcal H$ and put
\[
 \widetilde Q=\frac1{60}
   \sum_{r,s}T_{rs}\Sigma_r^{\mathrm{tot}}\Sigma_s^{\mathrm{tot}}.
\]
Compressing the inner derivative generator to $G$ gives
$I_s^{(j)}=\Sigma_s^G/5$: equivariance fixes the spin-one family up to a
scalar, and the highest weight gives $(1/2)/(5/2)=1/5$. At a point where
$T$ is diagonal, expansion of the total square therefore gives
\[
 T_D=\frac1{10}\sum_s t_s I_s^{(i)}\Sigma_s^G
     =3\widetilde Q-2Q_{5/2}(T),\qquad
 B_F(T)=Q_{5/2}(T)
\]
on the compressed inner factor. The inner square contributes
$40Q_{5/2}(T)$, while the trace-free outer half-spin square is zero.
Consequently
\[
 T_D+2B_F(T)=3\widetilde Q.
\]
This combined identity cancels the inner curvature action. On the free
$(3,1)$ type, $\widetilde Q=Q_3(T)\le4x/5$ and $U_D\le2y$. Thus the choice
\[
 \mathcal P=4-\frac56\bigl(T_D+2B_F(T)\bigr)-U_D
            =4-\frac52\widetilde Q-U_D
\]
has its plus and minus losses bounded by $2x$ and $2y$, respectively.
The scalar $4$ then meets the sectional boundary $x+y=2$. This explains
the useful coefficients; no uniqueness of the choice is asserted.

\begin{lemma}[Nonnegativity on the highest Hessian type]
\label{lem:hessian-operator}
The combined compressed operator $\mathcal P$ preserves all four types
in $\mathcal H$. On $(3,1)$ it satisfies
$\mathcal P\ge(4-2x-2y)\operatorname{Id}\ge0$, and on the trace
type it satisfies
\[
 \mathcal P\operatorname{Tr}^*
   =\operatorname{Tr}^*(4-B_F(T)).
\]
The individual derivative-slot and matrix-slot summands need not preserve
the four types.
\end{lemma}

\begin{proof}
The total plus generators preserve each joint type. Similarly,
$\sum_su_s(J_s^{(i)}+J_s^{(j)})^2=4U_D$, so $U_D$ vanishes on minus
spin zero and is $Q_1(U)$ on minus spin one. Its latter spectrum is the
spectrum of $U$. The displayed combined formula for $\mathcal P$ therefore
preserves the joint types. On $(3,1)$,
\[
 \mathcal P\ge4-2x-2y\ge0.
\]
On $(2,0)$ the denominator in $Q_2$ is $24$, so under the trace
identification $\widetilde Q=(2/5)Q_2(T)=(2/5)B_F(T)$ and $U_D=0$.
This proves the stated trace action.
\end{proof}

\subsection{The weighted square and its first-derivative form}

The undetermined Hessian component $Y$ is now isolated, and its
curvature-weighted quadratic form is nonnegative. We complete its weighted
square and integrate by parts to replace the Hessian by first derivatives.
The prescribed trace component must be subtracted in full.

Set
\[
 B=\frac{\langle B_F(T)Z,Z\rangle}{6},\qquad
 C_N=\frac{\langle B_F(N)Z,Z\rangle}{6},\qquad
 E_B=\frac{|Q_{5/2}(T)Z|^2}{6},
 \qquad N=(T^2)_0.
\]
For a scalar weight \(m\), write \(L_m=dm\otimes Z\), with no
symmetrization, and let \((L_m)_Y\) denote its \((3,1)\) projection.

\begin{proposition}[Weighted Hessian inequality]
\label{prop:weighted-hessian}
For every global Lipschitz scalar weight \mbox{$m\ge m_0>0$},
\begin{equation}
 0\le\mathbf E\left[
 mG_v(p,k,l)
 -\frac56\sum_i m_i\,
    \frac{\operatorname{tr}\bigl(Z_i\sum_jZ_j^2\bigr)}6
 +\frac{\langle\mathcal P(L_m)_Y,(L_m)_Y\rangle}{24m}
 \right],
 \tag{10}\label{eq:10}
\end{equation}
where \(i,j\) in the middle term are derivative indices and
\[
 G_v=\left(-42-3b^2-\frac{10}{3}v^2\right)e
       +75B+\frac{25}{3}C_N-\frac{100}{3}E_B
       +5v^2|p|^2.
\]
More precisely, if $\mathcal W_v(m)$ denotes the integrand in
\eqref{eq:10}, then the exact identity is
\begin{equation}
 \mathbf E\mathcal W_v(m)
 =\mathbf E\left[\frac m6\left\langle
 \mathcal P\left(Y+\frac{(L_m)_Y}{2m}\right),
               Y+\frac{(L_m)_Y}{2m}\right\rangle\right].
 \tag{10a}\label{eq:weighted-hessian-square}
\end{equation}
\end{proposition}

\begin{proof}
By Lemma~\ref{lem:hessian-operator},
\[
 \mathbf E\left[
 \frac{m}{6}\left\langle
 \mathcal P\left(Y+\frac{(L_m)_Y}{2m}\right),
              Y+\frac{(L_m)_Y}{2m}\right\rangle\right]\ge0.
\]
The last term in its expansion is the last term of \eqref{eq:10}.
Since \(\mathcal P\) preserves the type splitting and
\(X=Y+D_0\), the other terms are
\[
 \frac16\mathbf E\bigl[
 m\langle\mathcal PX,X\rangle+\langle\mathcal PX,L_m\rangle
 -m\langle\mathcal PD_0,D_0\rangle-\langle\mathcal PD_0,L_m\rangle
 \bigr].
\]
We compute these terms for smooth \(m\) first.  Compression may be omitted
inside any of the displayed inner products, because both arguments
belong to \(\mathcal H\).

\paragraph{The identity term.}
Integration by parts and \eqref{eq:1} give
\[
 \frac16\mathbf E\left[m|\nabla Z|^2
           +\sum_i m_i\langle\nabla_iZ,Z\rangle\right]
 =-\frac16\mathbf E\bigl[m\langle\Delta Z,Z\rangle\bigr]
 =\mathbf E\bigl[m(-13e+20B)\bigr].
\]
This will be multiplied by the coefficient \(4\) in \(\mathcal P\).

\paragraph{The derivative-slot curvature terms.}
Write \(A_{ikjl}=(T_D)_{ikjl}\).  The two terms to integrate are
\[
 \mathbf E\left[
 mA_{ikjl}\langle\nabla_iZ_j,\nabla_kZ_l\rangle
 +A_{ikjl}\langle\nabla_iZ_j,m_kZ_l\rangle\right].
\]
Integrate the derivative \(k\) on the second paired tensor.
The derivative of \(m\) cancels the second term exactly.
Moreover \(\nabla_kA_{ikjl}=0\), by divergence freedom and skewness in
\(i,k\).  The result is
\[
 -\mathbf E\bigl[mA_{ikjl}\langle\nabla_k\nabla_iZ_j,Z_l\rangle\bigr]
 =\frac12\mathbf E\bigl[
 mA_{ikjl}\langle[\nabla_i,\nabla_k]Z_j,Z_l\rangle\bigr].
\]
This calculation uses neither symmetry of \(X\) nor parallelness of
\(T_D\), and the factor \(1/2\) comes only from skewing the two
covariant derivatives.

The curvature in the derivative slot has the contraction
\[
 \frac12(T_D)_{ikjl}R_{ikjh}=3v^2\delta_{lh}.
\]
To see its normalization, substitute the quaternion expansions: the
contraction over \(i,k\) removes all mixed plus/minus terms and gives
\(\frac12\sum_s t_s(1+t_s)\delta_{lh}=3v^2\delta_{lh}\).
The matrix-slot curvature gives the quadratic form of
\[
 -\frac12\sum_s t_s(1+t_s)I_s\Sigma_s.
\]
Since \(t_s^2=N_{ss}+2v^2\), its compression is
\[
 -\frac12\bigl(8Q_{5/2}(T)+8Q_{5/2}(N)+2v^2C\bigr)
 =4v^2-4Q_{5/2}(T)-4Q_{5/2}(N).
\]
After division by \(6\), the full \(T_D\) contribution is therefore
\[
 \mathbf E\bigl[m(7v^2e-4B-4C_N)\bigr].
\]
For \(U_D\), the same integration gives \(3\mathbf E[mb^2e]\).
Its matrix-slot curvature contraction is zero, because the minus forms
are orthogonal to the plus forms defining the curvature action on
\(F\).

\paragraph{The nonparallel matrix-slot coefficient.}
Put \(A=B_F(T)\) temporarily.  One integration by parts gives
\[
 \mathbf E\left[m\langle A\nabla_iZ,\nabla_iZ\rangle
              +m_i\langle A\nabla_iZ,Z\rangle\right]
 =-\mathbf E\bigl[m\langle A\Delta Z,Z\rangle\bigr]
  -\mathbf E\bigl[m\langle(\nabla_iA)\nabla_iZ,Z\rangle\bigr].
\]
The operator \(\nabla_iA\) is symmetric, so the last term is
\[
 \frac12\mathbf E\left[
 \left\langle\left(m\Delta A+\sum_i m_i\nabla_iA\right)Z,Z\right\rangle
 \right].
\]
Here \(\nabla_iA=B_F(Z_i)\) and
\(\Delta A=6B_F(T)-6B_F(N)\), since \(B_F\) is a linear parallel
construction.  Also, pairing \(B_F(T)Q_{5/2}(T)Z\) with \(Z\) gives
\(|Q_{5/2}(T)Z|^2\), because the latter vector belongs to
\(G\).  After division by \(6\), the contribution is
\[
 \mathbf E\left[
 m(-10B-3C_N+20E_B)
 +\frac1{12}\sum_i m_i\langle B_F(Z_i)Z,Z\rangle\right].
\]
Since
\(\langle B_F(Z_i)Z,Z\rangle
  =\operatorname{tr}(Z_i\sum_jZ_j^2)\),
multiplication by \(-5/3\) produces exactly the middle term in
\eqref{eq:10}.  In particular, the variation of \(B_F(T)\) has been
retained explicitly; \(\mathcal P\) is not assumed parallel.

\paragraph{The prescribed trace component.}
By Lemmas~\ref{lem:hessian-components} and \ref{lem:hessian-operator},
\[
 \mathcal PD_0=\frac5{12}\operatorname{Tr}^*(4-B_F(T))\Delta T,
 \qquad \operatorname{Tr}L_m=\sum_i m_iZ_i.
\]
The characteristic polynomial of a trace-free \(3\times3\) matrix gives
\(B_F(T)T=N\) and \(B_F(T)N=v^2T\).  Thus
\[
 B_0=(4-B_F(T))(T-N)=(4+v^2)T-5N.
\]
Using \(\Delta T=6(T-N)\), the two trace subtractions, before division
by \(6\), are
\[
 -\frac52\mathbf E\left[
 m\langle B_0,\Delta T\rangle+\sum_i m_i\langle B_0,Z_i\rangle\right]
 =\frac52\mathbf E\bigl[m\langle\nabla B_0,Z\rangle\bigr].
\]
Since \(\nabla N=2B_F(T)Z\) and
\(\langle T,Z_i\rangle=6vp_i\), their normalized contribution is
\[
 \mathbf E\left[
 m\left(\left(10+\frac52v^2\right)e-25B+5v^2|p|^2\right)\right].
\]
Multiplying the preceding contributions by their coefficients in
\(\mathcal P\) and collecting gives \(G_v\) as stated.

Finally, a global Lipschitz weight $m\ge m_0>0$ can be approximated
uniformly and in $H^1$ by smooth positive weights with a common positive
lower bound. The first derivatives converge in $L^2$, so their quadratic
products in the square cost converge in $L^1$; the uniformly bounded
reciprocals of the weights pass to the same limit. All other terms use
only the weight and its first derivative. Hence the identity and inequality
extend to every weight in the stated class.
\end{proof}

The next algebraic expansion, as in \cite[Lemma~5.4]{OpenAI2026P},
expresses this invariant inequality in the gradient coordinates needed
for the coupling estimate.

\begin{lemma}[The quadratic-form matrices]
\label{lem:hessian-matrices}
At a point where \(v>0\), use the three vector parameters \(p,k,l\) in
\eqref{eq:2}.  Then \(G_v\) is the quadratic form, with scalar entries
multiplying vector dot products, whose matrix is
\begin{equation}
 \begin{pmatrix}
 -70-5b^2+100\sigma v-25v^2
   &-100\sqrt3\rho v&-35\rho v^2/\sqrt3\\
 -100\sqrt3\rho v
   &-210-15b^2-300\sigma v-190v^2&-75v-35\sigma v^2\\
 -35\rho v^2/\sqrt3
   &-75v-35\sigma v^2&-42-3b^2-10v^2
 \end{pmatrix}.
 \tag{11}\label{eq:11}
\end{equation}
\end{lemma}

\begin{proof}
Expanding the diagonal and off-diagonal matrix entries in \eqref{eq:2}
gives
\[
 \frac{B}{v}\ \longleftrightarrow\
 \begin{pmatrix}
 4\sigma/3&-4\rho/\sqrt3&0\\
 -4\rho/\sqrt3&-4\sigma&-1\\
 0&-1&0
 \end{pmatrix},
 \qquad
 \frac{C_N}{v^2}\ \longleftrightarrow\
 \begin{pmatrix}
 4/3&0&-\rho/\sqrt3\\
 0&-4&-\sigma\\
 -\rho/\sqrt3&-\sigma&0
 \end{pmatrix}.
\]
Here one uses
\(\sum s_i\widetilde s_i=0\),
\(\sum\widetilde s_i^3=-6\rho\), and
\(N/v^2=\operatorname{diag}(\sigma s_i+\rho\widetilde s_i)\).
The energy metric in these coordinates is
\(H=\operatorname{diag}(5/3,5,1)\).
If \(\mathsf B\) denotes the matrix of \(B\), the matrix of \(E_B\) is
\(\mathsf B H^{-1}\mathsf B\): \(B_F(T)\) compressed to
\(G\) is the self-adjoint operator whose quadratic form is
\(B\).  Using \(\sigma^2+\rho^2=1\), this gives
\[
 \frac{E_B}{v^2}\ \longleftrightarrow\
 \begin{pmatrix}
 16/15&0&4\rho/(5\sqrt3)\\
 0&21/5&4\sigma/5\\
 4\rho/(5\sqrt3)&4\sigma/5&1/5
 \end{pmatrix}.
\]
Substitution into the definition of \(G_v\) gives \eqref{eq:11} by
matrix addition.  The invariant expression for \(G_v\) is defined
everywhere; on the zero set of \(T\), all its derivative parameters
vanish almost everywhere, as established for the norm identity.
\end{proof}

\subsection{The two costs of a variable weight}

For the remaining argument choose the fixed positive weights
\begin{equation}
 M=1-\frac38(v-b),\qquad m=M^2,
 \qquad
 \overline M=1-\frac38(b-v),\qquad \overline m=\overline M^2.
 \tag{12}\label{eq:12}
\end{equation}
They are one rational choice for the later coupling argument. Since
$v+b\le2$, both $M$ and $\overline M$ lie in $[1/4,7/4]$;
in particular $m,\overline m\ge1/16$.
Almost everywhere,
\[
 dm=-\frac34M(p-q),\qquad d\overline m=-\frac34\overline M(q-p).
\]
Proposition~\ref{prop:weighted-hessian} applies to these global Lipschitz
weights. Their derivatives produce the two costs estimated next.

The Hessian has been eliminated from the integrated inequality. Two
terms caused by the weight remain: a cubic contraction multiplied by
$dm$, and a projected quadratic form in $dm\otimes Z$. We now bound
both by first-derivative energies, while keeping the sign needed for
the later scalar relaxation. The estimates and projection identities are
from \cite[Lemmas~5.5--5.6 and Proposition~5.7]{OpenAI2026P}; their
proofs below apply directly when $x+y\le2$.

\begin{lemma}[The cubic contraction]
\label{lem:cubic-contraction}
For every real tangent vector \(a\), with
\(Z_a=\sum_i a_iZ_i\),
\begin{equation}
 \left|\frac{\operatorname{tr}\bigl(Z_a\sum_jZ_j^2\bigr)}6\right|
 \le \frac45\sqrt{\frac35}\,|a|e^{3/2}.
 \tag{13}\label{eq:13}
\end{equation}
\end{lemma}

\begin{proof}
By homogeneity, assume \(|a|=1\); the case \(Z_a=0\) is immediate.
Put \(\nu=|Z_a|/\sqrt6\), and diagonalize the normalized matrix
\(s=Z_a/\nu\).  Apply the algebraic parametrization \eqref{eq:2} with
this \(s\), rather than with \(T/v\), and denote its vector parameters
by \(\widehat p,\widehat k,\widehat l\).
The diagonal entries of \(Z_a=\nu s\) imply
\(\widehat p_a=\nu\) and \(\widehat k_a=0\).
The part corresponding to the single component
\(\widehat p_a a\) has energy
\(e_1=5\nu^2/3\).  It is orthogonal in the energy metric to the
remaining parameters.  In the displayed matrix for \(B/v\), it has no
cross term with those parameters, because \(\widehat k_a=0\) and the
\(\widehat p,\widehat l\) entry is zero.  Its contribution to
\(\frac16\sum_j\operatorname{tr}(sZ_j^2)\) has absolute value at most
\(4e_1/5\).

For any symmetric trace-free \(3\times3\) matrix \(S\), its
characteristic polynomial gives
\[
 |(S^2)_0|^2=\frac{|S|^4}{6}.
\]
Since \(|s|=\sqrt6\), Cauchy--Schwarz yields
\(|\operatorname{tr}(sS^2)|\le |S|^2\).
Applied to the remaining tensor, this bounds its contribution by
\(e-e_1\).  Therefore
\[
 \left|\frac{\operatorname{tr}(Z_a\sum_jZ_j^2)}6\right|
 \le\nu\left(e-\frac{e_1}{5}\right)
 =\sqrt{\frac35}\,e^{3/2}
       \sqrt{\frac{e_1}{e}}\left(1-\frac{e_1}{5e}\right).
\]
For \(0\le r\le1\), the function \(\sqrt r(1-r/5)\) is increasing
and has maximum \(4/5\).  This proves \eqref{eq:13}.
\end{proof}

\begin{lemma}[Real-vector projection identities]
\label{lem:real-vector-projections}
Let \(P_+\) be the projection in \(\mathcal H\) onto total plus spin
\(3\), retaining both possible minus spins.  For any real tangent vector
\(a\) and any \(Z\in G\),
\[
 |P_+(a\otimes Z)|^2=\frac7{12}|a|^2|Z|^2,
\]
\[
 \left\langle\widetilde QP_+(a\otimes Z),P_+(a\otimes Z)\right\rangle
 =\frac12|a|^2\langle Q_{5/2}(T)Z,Z\rangle.
\]
\end{lemma}

\begin{proof}
Complexify at the point, writing the outer tangent representation as
\(A\otimes B\), with \(A,B\) the plus and minus spin halves, and
\(G_{\mathbb C}=H_{5/2}\otimes B'\).
For a real vector $a$, define the operator $R_a$ on $A$ by
$\operatorname{tr}(R_aD)=\langle a,(D\otimes\operatorname{Id}_B)a\rangle$
for every operator $D$ on $A$. This is the partial trace over $B$ of the
rank-one operator $a a^*$. We claim that
$R_a=\frac12|a|^2\operatorname{Id}_A$. Indeed, the three operators $-iI_s$ form a basis of the traceless
Hermitian operators on $A$, and each has zero expectation because $a^{\mathsf T}I_sa=0$.
The trace of $R_a$ is $|a|^2$, which fixes the scalar part.

The complete vector is \(a\otimes Z\), a product between the outer
factor and the entire inner factor.  Consequently, for every operator
\(O\) acting only on \(A\otimes H_{5/2}\),
\[
 \langle a\otimes Z,O(a\otimes Z)\rangle
 =\frac{|a|^2}{2}
    \langle Z,((\operatorname{Tr}_A O)\otimes\mathrm{Id}_{B'})Z\rangle.
\]
This partial trace retains every correlation of the inner plus factor
with \(B'\); neither \(a\) nor \(Z\) is assumed to be a product within
its own plus and minus factors.

The operator \(\operatorname{Tr}_A P_+\) is scalar on the irreducible
space \(H_{5/2}\).  Its trace is the dimension \(7\) of spin \(3\),
whereas \(H_{5/2}\) has dimension \(6\).  Thus it equals
\(7\mathrm{Id}/6\), proving the first identity.

The family
\(\operatorname{Tr}_A(P_+\widetilde Q(T)P_+)\) is equivariant and
linear in the spin-\(2\) variable \(T\).  Since spin \(2\) occurs once
in \(\operatorname{End}(H_{5/2})\), it is a scalar multiple of
\(Q_{5/2}(T)\).  To determine the scalar, take
\(T=\operatorname{diag}(-1,-1,2)\) and the inner highest weight
\(5/2\).  The outer weight \(1/2\) gives total weight \(3\) and
\(Q_3(T)\)-eigenvalue \(-1\).  The outer weight \(-1/2\) projects to
total weight \(2\), where that eigenvalue is
\((12-3\cdot2^2)/15=0\).  The normalized outer trace is therefore
\(-1/2\).  Since \(Q_{5/2}(T)\) has eigenvalue \(-1\) on its highest
weight, the second identity has coefficient \(1/2\), as asserted.
\end{proof}

\begin{proposition}[The projection cost]
\label{prop:hessian-projection-cost}
For the weight \eqref{eq:12},
\begin{equation}
 \frac{\langle\mathcal P(L_m)_Y,(L_m)_Y\rangle}{24m}
 \le\frac9{64}|p-q|^2
       \left[\frac7{12}(4+y)e-\frac54B\right].
 \tag{14}\label{eq:14}
\end{equation}
The bracket on the right is nonnegative.  Both assertions remain valid
under the closed inequalities \(v+b\le2\) and \(x+y\le2\).
\end{proposition}

\begin{proof}
Let \(P_-\) be the projection onto total minus spin \(1\), so the
highest Hessian projection is \(P_+P_-\).  On that type,
\(-U_D\le y\), and hence \(\mathcal P\le A\), where
\[
 A=4+y-\frac52\widetilde Q.
\]
This operator acts only on the plus factors and commutes with \(P_-\).
On the entire high-plus subspace,
\[
 A\ge4+y-2x\ge0;
\]
the last inequality follows from \(x+y\le2\).
It follows in operator order that \(P_-AP_-\le A\) there.
Thus, for every real \(a\),
\[
 \langle\mathcal P(a\otimes Z)_Y,(a\otimes Z)_Y\rangle
 \le \langle AP_+(a\otimes Z),P_+(a\otimes Z)\rangle
 =6|a|^2\left[\frac7{12}(4+y)e-\frac54B\right],
\]
where the last equality is Lemma~\ref{lem:real-vector-projections}.
This operator argument is valid for arbitrary correlations after the
plus projection and does not discard a mixed minus-factor term.

Almost everywhere,
\(dm=-\frac34M(p-q)\).
Substituting \(a=dm\), dividing by \(24m\), and using \(m=M^2\)
gives \eqref{eq:14}.
Finally, for any unit real \(a\), the same equality and the lower bound
for \(A\) give
\[
 \frac7{12}(4+y)e-\frac54B
 \ge \frac7{12}(4+y-2x)e\ge0.
\]
In particular, subsequent upper estimates may replace \(|p-q|^2\) by
\((|p|+|q|)^2\) in \eqref{eq:14} without reversing the inequality.
\end{proof}
